% ===== Abstract =====
\begin{center}
{\color{deeprose}\bfseries\large Abstract}
\end{center}

\vspace{0.4em}

\noindent
This paper reframes the construction of a generative relational field as the proper object of an aesthetic education of the everyday. It proceeds by a reduction chain: relational production requires a field, a conditional environment of generative coupling; the construction of such a field is an aesthetic rather than an engineering matter, because the appropriateness it must achieve is state-dependent and cannot be specified by any rule; the aesthetic and the ethical are coupled at their root, since the perception of the fitting and the perception of the good are one perception of a field's generative direction; and therefore the cultivation of the capacity to build generative fields is what aesthetic education most fundamentally names. The paper establishes a concept, \emph{relational aesthetics}, on which beauty is a coupling event between subjects rather than a property of objects or a state of subjects, and a four-step model of cultivation, \emph{perception}, \emph{reception}, \emph{sharing}, and \emph{reproduction}, shown to be the smallest form a good cycle takes when enacted between two people, a holonomy cycle that returns to its root and not its origin. It reconstructs the dialectic of principal and secondary contradiction into a generative form that directs attention rather than concentrating force, in keeping with a dialectical-materialist commitment that treats no single theory as a final rationality but assigns each its determinate place and limit. At four points where the model's movements carry the densest theory, the paper sets the contending traditions of predictive processing, Buddhist and Daoist contemplative analysis, phenomenology, psychoanalysis, the ethics of the other, and the dialogical and sociological accounts of the shared into adversarial comparison, resolving them not by a synthesis that absorbs them but by a layered arrangement that assigns each its level, marks what it can and cannot explain, and keeps the differences while finding the common ground. An orthogonality guard separates the refinement of the sensibility from the justice of the field, against aestheticism; an implementation-level interface specifies the observable signatures each movement should leave, without reducing the relational to the neural; and a political economy of the field argues that the cultivation, requiring no surplus of time or means, is available precisely to those whose days hold no surplus, and is most needed by them. The everyday, perceived and received and shared and renewed with a trained and unattached attention, is shown to be where the good cycle is most needed and most possible, and its cultivation the aesthetic education that the whole argument was for.

\vspace{1em}

\noindent{\color{deeprose}\bfseries Keywords:} \enspace relational aesthetics; aesthetic education; holonomy and the good cycle; the construction of fields; generative dialectical materialism; the everyday.
